\documentclass[t,
handout,
]{beamer}

%beamer settings
\setbeamerfont{normal text}{size=\large}
\AtBeginDocument{\usebeamerfont{normal text}}

% general language
\usepackage{polyglossia}
\setmainfont{Arial}
\setdefaultlanguage{english}
%\usepackage[utf8]{inputenc} % ignor1 anyway due to uft8 based engines
\usepackage[autostyle=true]{csquotes}
\setcounter{tocdepth}{3}
\renewcommand{\baselinestretch}{1.0}

%colors
\usepackage{xcolor}
\definecolor{bluei}{rgb}{0.0, 0.0, 1.0}
\definecolor{greni}{rgb}{0.3, 0.8, 0.2}
\definecolor{greyi}{rgb}{0.5, 0.5 0.5}
\definecolor{b1}{rgb}{0, 0, 0}
\definecolor{r1}{rgb}{255, 0, 0}

% graphics
\usepackage{amssymb}
\usepackage{graphicx}
\usepackage{subfigure}
\usepackage{subcaption}

%date
\usepackage[nodayofweek,level]{datetime}

%change
\newcommand{\authorName}{Janis Jacobi\\Niels Schröter}
\newcommand{\topicName}{physics modelling agent}
\newdate{date1}{24}{12}{2025}

% general layout settings for beamer
\usetheme{Madrid}
%\usecolortheme[RGB={0, 100, 250}]{beamer color theme}
\usecolortheme{seahorse}
%\usefonttheme{structurebold}

\useoutertheme{default}
% beamer template
\setbeamercolor{footline left}{bg=blue!20, fg=b1}
\setbeamercolor{footline center}{bg=blue!10, fg=b1}
\setbeamercolor{footline right}{bg=blue!20, fg=b1}
\setbeamertemplate{footline}{%
	\leavevmode%
	\hbox{%
		\begin{beamercolorbox}[wd=.33\paperwidth,ht=2.5ex,dp=1ex,left]{footline left}%
			\hspace{1em}\topicName
		\end{beamercolorbox}%
		\begin{beamercolorbox}[wd=.34\paperwidth,ht=2.5ex,dp=1ex,center]{footline center}%
			\insertshortdate{}
		\end{beamercolorbox}%
		\begin{beamercolorbox}[wd=.33\paperwidth,ht=2.5ex,dp=1ex,right]{footline right}%
			\insertframenumber{} / \inserttotalframenumber\hspace{1em}
		\end{beamercolorbox}%
	}%
}

\setbeamertemplate{headline}{}
% restore visible bullets



% title-page settings
\title{\topicName}
\subtitle{update, Arxiv and units}
\date{\today}
%\logo{logo text}
\author{\authorName}

\usepackage[backend=biber, style=numeric, citestyle=numeric, hyperref=true, date=short, backref=false]{biblatex}
%\bibliography{biblio/jabRef.bib}
%\addbibresource{/biblio/jabRef.bib}
\setlength\bibitemsep{7pt}
\setcounter{biburllcpenalty}{9000}

%captions
\usepackage{caption}
\usepackage{subcaption}

%math
\usepackage{mathtools}
\usepackage{unicode-math}
\setmathfont{Latin Modern Math}
\usepackage{amsmath,amssymb}
\usepackage[locale=DE, separate-uncertainty=true]{siunitx}
\usepackage{nicefrac}
\usepackage{physics}
\newcommand{\RNum}[1]{\uppercase\expandafter{\romannumeral #1\relax}} %roman numeral
\numberwithin{equation}{section}
\usepackage{numprint}
\usepackage[version=4]{mhchem}

\newcommand{\al}{\ensuremath{\mathcolor{c1}{|}}}
\newcommand{\scProd}[2]{\ensuremath{\langle #1,#2 \rangle}}
\newcommand{\parti}[2]{\ensuremath{\frac{\mathrm{d}#1}{\mathrm{d}#2}}}
\newcommand{\pard}[2]{\ensuremath{\frac{\partial #1}{\partial #2}}}

% graphics
\usepackage{graphicx}
%\usepackage[margin=2cm]{geometry}
\usepackage{subcaption}
\captionsetup{compatibility=false}

%circuits and drawing
\usepackage{tikz}
\usetikzlibrary{
	shapes.geometric,
	arrows.meta,
	positioning,
	fit,
	calc
}

% insert documents
\usepackage{pdfpages}
\usepackage{float}

% tables
\usepackage{tabularx}
\usepackage{booktabs}
\usepackage{longtable}
\usepackage{pgfplotstable}
\usepackage{multirow}

%csv
\usepackage{csvsimple}

%code
\usepackage{listings}
\usepackage{lstautogobble}
\lstdefinestyle{py}{language={Python}}
\lstset{style=py,
	basicstyle=\normalsize,
	frame=tb,
	language=Python,
	breaklines=true,
	showstringspaces=false,
	columns=flexible,
	stringstyle = \color{blue},
	commentstyle = \color{grey},
	tabsize=3,
	numbers=left,
	keywordstyle=\color{r1},
	columns=flexible,
	frame=single,
	xleftmargin = 20pt,
	autogobble=true,
}

% beamer content display settings:
\usepackage{overpic}

%BIBLIOGRAPHIE
\usepackage[
	backend=biber, %type of file
	style=numeric, %bib entry listing
	citestyle=numeric, % how citation is in text
	hyperref=true, %link the bib entryies
%	date=short, %date format
%	backref=true, % where citation was used
%	url=true, %display url
%	doi=true, % digital object identifier: a permanent, unique link
%	maxbibnames=10, %max number of names displayed
%	minbibnames=3, %minimal number of names displayed
]{biblatex}

%\addbibresource{bib/bibliography.bib}
\setlength\bibitemsep{7pt} %spacing between entries
\setcounter{biburlucpenalty}{9000} %prevent mid word url breaks
\setcounter{biburllcpenalty}{9000} %makes sure urls do not run out of the page

%hyperlinks
\usepackage{hyperref}
\hypersetup{
	baseurl={https://www.janisjacobi.de/home.html},
	pdfpagemode = UseNone,
	pdfauthor={\authorName},
	pdfsubject = {\topicName},
	pdfcreator = {\authorName},
	pdfcreationdate = {\today},
	pdfnewwindow,
	bookmarksopen=true,
	backref=page,
	colorlinks,
	urlcolor=bluei,
	citebordercolor=bluei,
	citecolor=bluei,
	linkcolor=bluei,
}

\begin{document}
	\begin{frame}
		\titlepage
	\end{frame}
	
	\begin{frame}{outline}
		\tableofcontents
	\end{frame}
	
	\section{arxiv problem}
	\begin{frame}{the arxiv issue}[c]
		\begin{alertblock}{the problem}
			\begin{itemize}
				\item arxiv queries where rejected
				\item papers could not be read or downloaded
			\end{itemize}
		\end{alertblock}
		\begin{block}{the underlying issue}
			\begin{itemize}
				\item crewai's ArxivPaperTool does not implement a proper exponential back-off and waiting time for the arxiv request
				\item[$\implies$] Arxiv-API-server rejects requests and even blocks IP address temporary
			\end{itemize}
		\end{block}
	\end{frame}
		
	\begin{frame}{solution}
		\begin{exampleblock}{solution}
			\begin{itemize}
				\item exponential back-off and proper waiting time (\si{5 \second} between request)
				\item split the crew in two:
				\begin{enumerate}
					\item arxiv request and paper downloading (all in one step)
					\item extract infos from downloaded papers; build the model and 	implement it
				\end{enumerate}
			\end{itemize}
		\end{exampleblock}
		
		\begin{block}{further improvements}
			\begin{itemize}
				\item from one set of arxiv papers, the second crew can be run in multiple configurations without a new arxiv-server request
				\item[$\implies$] drastically reduces the number of server requests
				\item[$\implies$] papers could be added from another source into the directory the second crew reads from
			\end{itemize}
		\end{block}
	\end{frame}
	
	\section{units}
	
	\begin{frame}{the key formula}
		with
		{\centering
		\begin{equation*}
			M_y = \frac{\mu_b \abs{e} \tau}{2\pi} m^* \alpha_R \left[\hat{z} \times \vec{E}\right]_y
			\end{equation*}}
		
		\begin{itemize}
			\item $\mu_B$ Bohr magneton in \si{\joule /\tesla}
			\item $e$ charge in \si{\coulomb}
			\item $\tau$ relaxation time in \si{\second}
			\item $m^*$ effectice mass in \si{\kilo \gram}
			\item $\alpha_R$ rashba coupling strength in (\si{\joule \meter})
			\item $\vec{E}$ electric field in \si{\volt / \meter } $=$\si{\kilogram \meter / (\ampere \second^3)}
			\item $M_y$ magnetization expected in \si{\ampere/ \meter}
		\end{itemize}
		 
		changing $\alpha_R$ units to \si{\meter^4 /\second} would result in matching units
	\end{frame}

	
%	\begin{frame}[allowframebreaks]{sources}
%		\printbibliography
%	\end{frame}
%	
%	\begin{frame}[t]{AI usage}
%		The drawings were created with help of \href{https://copilot.microsoft.com/chats}{microsoft copilot}. I take full responsibility for the figures contents.
%	\end{frame}
\end{document}